% Abhakseite „Das kann ich“ – Zone nur im Gesamt (\mitzone)
%% TODO Ich-kann-Sätze: Platzhalter wie in den Aufgabendateien
\begin{abhakseite}
\ifdefined\mitzone
\abhakgruppe{Kennst du schon}
\abhak{1}{Ergebnismengen zweistufiger Experimente aufzählen und günstige Fälle abzählen (Paare, Reihenfolge beachten)}
\abhak{2}{Brüche addieren und zur Summe eins ergänzen}
\abhak{3}{Säulendiagramme lesen (Höhen, Symmetrie, Verhältnisse)}
\abhak{4}{Ergebnismengen zweistufiger Experimente aufzählen und günstige Fälle abzählen (Paare, Reihenfolge beachten) – Fehler finden}
\abhak{5}{Ergebnismengen zweistufiger Experimente aufzählen und günstige Fälle abzählen (Paare, Reihenfolge beachten) – selbst rechnen}
\fi
\abhakgruppe{Einheit 1 · Verteilung aufstellen: die Werte der Zufallsgröße aus den Regeln gewinnen (alle Ergebnisfolgen durchrechnen), die Tabelle durch Abzählen füllen (günstige Paare je Wert), fehlende Wahrscheinlichkeiten über die Summe eins}
\abhak{6}{„Was muss eins ergeben?“}
\abhak{7}{Verteilung aufstellen}
\abhak{8}{Verteilung aufstellen – Fehler finden, Begründen, Anwendung, Darstellungswechsel}
\abhakgruppe{Einheit 2 · Einheit 2}
\abhak{9}{Verteilung lesen}
\abhak{10}{Verteilung lesen – Fehler finden, Begründen, Anwendung, Darstellungswechsel}
\end{abhakseite}
